\chapter{Ricerca e ordinamento}
\label{chap:ricerca-ordinamento}

La ricerca individua un elemento che soddisfa una condizione, mentre
l'ordinamento dispone gli elementi secondo una relazione prestabilita.
Questi problemi permettono di confrontare strategie con costi differenti,
utilizzando le notazioni della \cref{sec:fi-complessita-asintotica}.
Si considerano liste Python di interi, già introdotte nella
\cref{sec:fi-python-liste}, e si indica con \(n\) la loro lunghezza.
Nei programmi il parametro \texttt{v\_size} coincide con \texttt{len(v)}.
Si assume accesso per indice in tempo costante, come per gli array
della \cref{sec:fi-strutture-array}.\footnote{Materiale
di riferimento: S. Canazza, \emph{Ricerca e Ordinamento}, dispense del
corso di Fondamenti di Informatica.}

\section{Ricerca lineare}
\label{sec:fi-ricerca-lineare}

Data una lista \(v\) e un valore \(x\), la \emph{ricerca lineare}, o
\emph{sequenziale}, esamina gli elementi nell'ordine degli indici e
restituisce la posizione della prima occorrenza di \(x\). Se il valore
non compare, restituisce \(-1\). In questo contesto \(-1\) è un valore
convenzionale che segnala l'insuccesso, non un indice da utilizzare
per accedere alla lista.

\begin{unipdcode}[language=Python]
def linear_search(v, v_size, value):
    for i in range(v_size):
        if v[i] == value:
            return i
    return -1
\end{unipdcode}


Prima di esaminare l'indice \(i\), nessuno degli elementi nelle
posizioni da \(0\) a \(i-1\) è uguale a \(x\). Se il confronto riesce,
l'indice restituito è pertanto la prima occorrenza; in caso contrario,
l'algoritmo termina dopo aver esaminato l'intera lista.

Se il valore cercato occupa la posizione \(i\), si effettuano \(i+1\)
confronti; se è assente, se ne effettuano \(n\). Per \(n\ge1\) il caso
migliore ha costo \(\Theta(1)\), mentre quello peggiore ha costo
\(\Theta(n)\). Condizionatamente alla presenza del valore e assumendo
che la prima occorrenza sia uniformemente distribuita fra le \(n\)
posizioni, il numero medio di confronti è
\[
  \frac{1+2+\cdots+n}{n}=\frac{n+1}{2}\in\Theta(n).
\]
La ricerca richiede \(\Theta(1)\) memoria ausiliaria.

\section{Ricerca binaria}
\label{sec:fi-ricerca-binaria}

Se la lista è ordinata in senso non decrescente, la ricerca può
escludere a ogni confronto una parte degli elementi. La
\emph{ricerca binaria} confronta il valore cercato con quello centrale
dell'intervallo ancora da esaminare: se i valori differiscono,
prosegue soltanto nella metà compatibile con l'ordine della lista.

\clearpage

\begin{unipdcode}[language=Python]
def binary_search(v, v_size, value):
    return bin_search(v, 0, v_size - 1, value)


def bin_search(v, from_index, to_index, value):
    if from_index > to_index:
        return -1

    mid = (from_index + to_index) // 2
    if v[mid] == value:
        return mid
    if v[mid] < value:
        return bin_search(v, mid + 1, to_index, value)
    return bin_search(v, from_index, mid - 1, value)
\end{unipdcode}


La procedura ricorsiva mantiene un intervallo di indici
\([\ell,r]\). Se \(\ell>r\), l'intervallo è vuoto e il valore non è
presente. Altrimenti, posto \(m=\lfloor(\ell+r)/2\rfloor\),
si confronta \(x\) con \(v[m]\). Quando \(v[m]<x\), tutti gli elementi
di indice non superiore a \(m\) possono essere esclusi; quando
\(v[m]>x\), si possono escludere quelli di indice non inferiore a
\(m\). Questa proprietà, insieme alla terminazione su intervalli
sempre più piccoli, giustifica la correttezza dell'algoritmo.

Per esempio, nella lista \([5,9,11,12,17,21]\), la ricerca di
\(12\) confronta successivamente gli elementi di indici \(2\), \(4\)
e \(3\), restituendo quest'ultimo indice. In presenza di duplicati,
la ricerca binaria non garantisce di individuare la prima occorrenza.

Nel caso peggiore, dopo ogni confronto rimangono al più
\(\lfloor n/2\rfloor\) elementi. Per un costo unitario del confronto,
la ricorrenza si esprime, a meno dei termini costanti, come
\[
  T(n)=T(\lfloor n/2\rfloor)+\Theta(1).
\]
Dopo \(k\) riduzioni l'intervallo ha dimensione al più
\(\lfloor n/2^k\rfloor\); bastano quindi \(\Theta(\log n)\)
confronti nel caso peggiore. Il caso migliore richiede un solo
confronto. La versione ricorsiva usa \(O(\log n)\) spazio per la
pila delle chiamate, secondo il modello della
\cref{sec:fi-pila-ricorsione}. L'ordinamento iniziale è una
precondizione: se la lista non è ordinata, il procedimento può
restituire erroneamente \(-1\).

\section{Ricerca del massimo e del minimo}
\label{sec:fi-ricerca-estremi}

La ricerca di un estremo in una lista non ordinata richiede di
esaminare ogni elemento. Per una lista non vuota, si assume il primo
valore come candidato e lo si aggiorna quando se ne incontra uno
rispettivamente maggiore o minore.

\begin{unipdcode}[language=Python]
def find_max(v, v_size):
    if v_size == 0:
        raise ValueError("Lista vuota")
    massimo = v[0]
    for i in range(1, v_size):
        if v[i] > massimo:
            massimo = v[i]
    return massimo


def find_min(v, v_size):
    if v_size == 0:
        raise ValueError("Lista vuota")
    minimo = v[0]
    for i in range(1, v_size):
        if v[i] < minimo:
            minimo = v[i]
    return minimo
\end{unipdcode}


Dopo aver esaminato i primi \(i\) elementi, la variabile candidato
contiene l'estremo di quella parte della lista. Al termine il
risultato è quindi l'estremo dell'intera lista. In entrambi i casi
si effettuano esattamente \(n-1\) confronti per \(n\ge1\): la
complessità temporale è \(\Theta(n)\) e quella spaziale ausiliaria
\(\Theta(1)\). La lista vuota non ammette massimo né minimo, e le
funzioni segnalano tale condizione con \texttt{ValueError}.

\section{Ordinamento per selezione}
\label{sec:fi-ordinamento-selezione}

L'\emph{ordinamento per selezione} (\emph{Selection Sort}) dispone
gli elementi in ordine non decrescente. A ogni iterazione individua
il minimo nella parte non ancora ordinata e lo scambia con il primo
elemento di tale parte. La ricerca dell'estremo riprende il
procedimento della \cref{sec:fi-ricerca-estremi}, limitandolo a un
intervallo di indici.

\begin{unipdcode}[language=Python]
def selection_sort(v, v_size):
    for i in range(v_size - 1):
        indice_min = i
        for j in range(i + 1, v_size):
            if v[j] < v[indice_min]:
                indice_min = j
        if indice_min != i:
            v[i], v[indice_min] = v[indice_min], v[i]
\end{unipdcode}


All'inizio dell'iterazione di indice \(i\), i primi \(i\) elementi
sono ordinati e contengono i \(i\) valori minori della lista, tenendo
conto delle eventuali ripetizioni. La selezione del minimo nella
parte restante conserva tale proprietà per \(i+1\); quando rimane
un solo elemento, la lista è ordinata.

Per esempio, la lista \( [11,9,17,5,12] \) viene trasformata in
\( [5,9,11,12,17] \). Le ricerche del minimo richiedono
\[
  (n-1)+(n-2)+\cdots+1
  =\frac{n(n-1)}2\in\Theta(n^2)
\]
confronti, indipendentemente dalla disposizione iniziale. Gli scambi
sono al più \(n-1\). L'algoritmo modifica la lista originale,
richiede \(\Theta(1)\) spazio ausiliario e, nella forma mostrata,
non è stabile: lo scambio può invertire l'ordine relativo di due
elementi con chiave uguale.

La stessa strategia ammette una realizzazione ricorsiva: dopo aver
fissato il minimo, si ordina la parte rimanente. Il caso base è
una parte di lunghezza al più uno.

\clearpage

\begin{unipdcode}[language=Python]
def selection_sort_ricorsivo(v, inizio, v_size):
    if inizio >= v_size - 1:
        return
    indice_min = inizio
    for j in range(inizio + 1, v_size):
        if v[j] < v[indice_min]:
            indice_min = j
    if indice_min != inizio:
        v[inizio], v[indice_min] = v[indice_min], v[inizio]
    selection_sort_ricorsivo(v, inizio + 1, v_size)
\end{unipdcode}


Anche questa versione effettua \(\Theta(n^2)\) confronti; la
profondità delle chiamate cresce però come \(\Theta(n)\), anziché
richiedere memoria ausiliaria costante.

\section{Ordinamento per fusione}
\label{sec:fi-ordinamento-fusione}

L'\emph{ordinamento per fusione} (\emph{MergeSort}) applica la
strategia \emph{divide et impera}: divide la lista in due parti di
lunghezza pressoché uguale, ordina ricorsivamente ciascuna parte e
fonde i risultati. Una lista con meno di due elementi costituisce
il caso base. La suddivisione dell'ingresso
\([5,7,9,1,8]\) è illustrata nella
\cref{fig:fi-merge-sort-suddivisione}.

\begin{figure}[htbp]
  \centering
  \begin{tikzpicture}[unipd/diagram,x=1cm,y=1cm]
  \node[unipd/node] (a) at (0,0) {5, 7, 9, 1, 8};
  \node[unipd/node] (b) at (-3,-1.3) {5, 7, 9};
  \node[unipd/node] (c) at (3,-1.3) {1, 8};
  \node[unipd/node] (d) at (-4.1,-2.6) {5, 7};
  \node[unipd/node] (e) at (-1.9,-2.6) {9};
  \node[unipd/node] (f) at (2,-2.6) {1};
  \node[unipd/node] (g) at (4,-2.6) {8};
  \node[unipd/node] (h) at (-4.7,-3.9) {5};
  \node[unipd/node] (i) at (-3.5,-3.9) {7};
  \draw[unipd/directed edge] (a) -- (b);
  \draw[unipd/directed edge] (a) -- (c);
  \draw[unipd/directed edge] (b) -- (d);
  \draw[unipd/directed edge] (b) -- (e);
  \draw[unipd/directed edge] (c) -- (f);
  \draw[unipd/directed edge] (c) -- (g);
  \draw[unipd/directed edge] (d) -- (h);
  \draw[unipd/directed edge] (d) -- (i);
\end{tikzpicture}

  \caption{Suddivisione ricorsiva delle liste in MergeSort; la fusione
  delle parti ordinate avviene risalendo l'albero.}
  \label{fig:fi-merge-sort-suddivisione}
\end{figure}

La \emph{fusione} (\emph{merge}) di due liste già ordinate confronta
i rispettivi primi elementi ancora disponibili, copia il minore
nella lista di destinazione e avanza l'indice corrispondente. Quando
una parte è esaurita, copia gli elementi residui dell'altra.

\clearpage

\begin{unipdcode}[language=Python]
def merge_sort(v, v_size):
    if v_size < 2:
        return
    mid = v_size // 2
    left = v[:mid]
    right = v[mid:v_size]
    merge_sort(left, len(left))
    merge_sort(right, len(right))
    merge(v, left, right)


def merge(v, v1, v2):
    i = i1 = i2 = 0
    while i1 < len(v1) and i2 < len(v2):
        if v1[i1] <= v2[i2]:
            v[i] = v1[i1]
            i1 += 1
        else:
            v[i] = v2[i2]
            i2 += 1
        i += 1
    while i1 < len(v1):
        v[i] = v1[i1]
        i += 1
        i1 += 1
    while i2 < len(v2):
        v[i] = v2[i2]
        i += 1
        i2 += 1
\end{unipdcode}


Durante la fusione, la parte di destinazione già scritta è ordinata
e contiene i più piccoli fra gli elementi finora considerati. Poiché
le due parti sono ordinate, il prossimo elemento da copiare è uno
dei loro primi elementi disponibili. La proprietà è conservata a
ogni passo, e al termine tutti gli elementi sono stati copiati.
La condizione \lstinline|<=| nella fusione privilegia, a parità di
valore, l'elemento della parte sinistra: questa implementazione è
quindi \emph{stabile}. Per \(v=[5,2,4,1,3]\), la chiamata
\texttt{merge\_sort(v, len(v))} produce \(v=[1,2,3,4,5]\).

La creazione delle due sottoliste e la loro fusione richiedono
\(\Theta(n)\) operazioni complessive. Indicando con \(T(n)\) il costo
per una lista di \(n\) elementi, si ottiene
\[
  T(n)=T(\lfloor n/2\rfloor)
       +T(\lceil n/2\rceil)+\Theta(n),
  \qquad T(1)=\Theta(1).
\]
Per \(n=2^k\), se il costo non ricorsivo è \(cn\), la ricorrenza
si riduce a
\[
  T(n)=2T(n/2)+cn
      =2^kT(n/2^k)+kcn
      =nT(1)+cn\log_2n.
\]
Ne segue \(T(n)\in\Theta(n\log n)\); la stessa stima vale per
dimensioni non potenze di due, poiché l'albero delle chiamate ha
\(\Theta(\log n)\) livelli e ciascun livello richiede nel complesso
\(\Theta(n)\) lavoro. La base del logaritmo non influisce sulla
classe asintotica, perché un cambiamento di base introduce soltanto
un fattore costante. Il \emph{teorema maestro} fornisce un criterio
generale per alcune ricorrenze della forma
\(T(n)=aT(n/b)+f(n)\); in questo caso basta la sostituzione
precedente. La suddivisione mediante slicing e la fusione
richiedono \(\Theta(n)\) memoria ausiliaria complessiva al picco,
oltre alla pila di profondità \(\Theta(\log n)\).

\section{Ordinamento per inserimento}
\label{sec:fi-ordinamento-inserimento}

Nell'\emph{ordinamento per inserimento} (\emph{Insertion Sort}) la
parte iniziale della lista è mantenuta ordinata. Ogni nuovo elemento
viene conservato temporaneamente, mentre gli elementi maggiori
vengono spostati verso destra fino a liberare la posizione corretta.

\begin{unipdcode}[language=Python]
def insertion_sort(v, v_size):
    for i in range(1, v_size):
        temp = v[i]
        j = i
        while j > 0 and temp < v[j - 1]:
            v[j] = v[j - 1]
            j -= 1
        v[j] = temp
\end{unipdcode}


Prima dell'iterazione di indice \(i\), il prefisso di lunghezza \(i\)
è ordinato. Il ciclo interno sposta soltanto elementi maggiori del
valore da inserire; dopo l'inserimento, il prefisso di lunghezza
\(i+1\) risulta ancora ordinato. L'uso del confronto stretto
\lstinline|<| mantiene l'ordine relativo degli elementi uguali, rendendo
l'algoritmo stabile.

Quando la lista è già ordinata, il ciclo interno non compie
spostamenti: il costo del caso migliore è \(\Theta(n)\).
Se invece gli elementi sono in ordine strettamente decrescente,
l'iterazione \(i\) richiede \(i\) spostamenti e il numero totale è
\[
  \sum_{i=1}^{n-1}i=\frac{n(n-1)}2\in\Theta(n^2).
\]
Per una permutazione uniformemente casuale di valori distinti,
ciascun nuovo elemento supera in media metà dei precedenti; il
costo medio è pertanto \(\Theta(n^2)\). La memoria ausiliaria è
\(\Theta(1)\). In presenza di pochi elementi fuori posizione,
Insertion Sort può risultare vantaggioso rispetto a metodi con un
costo indipendente dall'ordine iniziale.

\section{Ordinamento rapido}
\label{sec:fi-ordinamento-rapido}

L'\emph{ordinamento rapido} (\emph{QuickSort}) è un altro algoritmo
\emph{divide et impera}. Scelto un elemento \emph{pivot}, la fase di
\emph{partizione} riordina il tratto considerato ponendo a sinistra
gli elementi minori o uguali al pivot e a destra quelli maggiori.
Il pivot viene quindi collocato nella sua posizione definitiva e il
procedimento prosegue ricorsivamente sui due intervalli rimanenti.

\begin{unipdcode}[language=Python]
def quicksort(a, inizio, fine):
    if inizio < fine:
        posizione_pivot = partiziona(a, inizio, fine)
        quicksort(a, inizio, posizione_pivot - 1)
        quicksort(a, posizione_pivot + 1, fine)

def partiziona(a, inizio, fine):
    pivot = a[fine]
    i = inizio - 1
    for j in range(inizio, fine):
        if a[j] <= pivot:
            i += 1
            a[i], a[j] = a[j], a[i]
    a[i + 1], a[fine] = a[fine], a[i + 1]
    return i + 1
\end{unipdcode}


Nella procedura di partizione mostrata, il pivot è l'ultimo elemento.
Durante la scansione, la parte del tratto già attraversata è divisa
fra valori non maggiori del pivot e valori maggiori. Scambiando
ciascun nuovo valore non maggiore con la prima posizione del secondo
gruppo, tale proprietà viene conservata; lo scambio conclusivo
colloca il pivot fra i due gruppi. Ogni chiamata ricorsiva opera su
un intervallo strettamente più piccolo, fino al caso base di
lunghezza al più uno.

Per esempio, per \(a=[5,2,9,1,7]\) la prima partizione pone
\(7\) nella posizione definitiva e lascia \(5,2,1\) a sinistra
e \(9\) a destra; al termine si ottiene \(a=[1,2,5,7,9]\).

La partizione di \(n\) elementi costa \(\Theta(n)\). Con divisioni
bilanciate il tempo soddisfa
\(T(n)=2T(n/2)+\Theta(n)\), quindi è
\(\Theta(n\log n)\). Nel caso peggiore, per esempio con una lista già
ordinata e scelta sistematica dell'ultimo elemento come pivot,
una delle parti è vuota e l'altra ha \(n-1\) elementi:
\[
  T(n)=T(n-1)+\Theta(n)\in\Theta(n^2).
\]
Con pivot scelto uniformemente a caso, il tempo \emph{atteso} è
\(\Theta(n\log n)\), senza richiedere che ogni partizione sia
bilanciata. La partizione opera direttamente sulla lista e usa
\(\Theta(1)\) spazio ausiliario; la pila ricorsiva occupa
\(O(\log n)\) spazio atteso e \(O(n)\) nel caso peggiore.
La versione mostrata non è stabile.

La scelta del pivot può utilizzare il primo elemento, l'ultimo,
quello centrale, un elemento casuale o la mediana di tre elementi.
Per migliorare il comportamento pratico si possono impiegare
Insertion Sort sulle parti molto piccole, partizioni a tre vie in
presenza di numerosi duplicati ed eliminazione delle chiamate
ricorsive in coda, richiamata nella
\cref{sec:fi-ricorsione-coda}. Le due partizioni indipendenti si
prestano inoltre all'esecuzione parallela.

\section{Confronto degli algoritmi}
\label{sec:fi-ordinamento-confronto}

Le complessità temporali dei quattro algoritmi studiati sono
riepilogate nella tabella seguente. Le stime si riferiscono a
confronti e accessi a elementi di dimensione costante, con la
convenzione introdotta nella \cref{sec:fi-complessita-modello}.

\begin{table}[htbp]
  \centering
  \caption{Complessità temporale degli algoritmi di ordinamento.}
  \label{tab:fi-confronto-ordinamenti}
  \begin{tabular}{lccc}
    \hline
    Algoritmo & Caso migliore & Caso medio & Caso peggiore\\
    \hline
    Selection Sort & \(\Theta(n^2)\) & \(\Theta(n^2)\) & \(\Theta(n^2)\)\\
    MergeSort & \(\Theta(n\log n)\) & \(\Theta(n\log n)\) & \(\Theta(n\log n)\)\\
    Insertion Sort & \(\Theta(n)\) & \(\Theta(n^2)\) & \(\Theta(n^2)\)\\
    QuickSort & \(\Theta(n\log n)\) & \(\Theta(n\log n)\) & \(\Theta(n^2)\)\\
    \hline
  \end{tabular}
\end{table}

Le stime medie assumono una distribuzione uniforme delle permutazioni
di valori distinti; per QuickSort il risultato riguarda anche la
versione con pivot casuale. MergeSort garantisce un costo
\(\Theta(n\log n)\) e conserva l'ordine degli elementi uguali,
ma richiede una lista ausiliaria. QuickSort limita le copie,
mentre Insertion Sort è particolarmente adatto a liste piccole o
quasi ordinate. La scelta dipende quindi non soltanto dalla
complessità asintotica, ma anche dalla disposizione dei dati,
dai requisiti di stabilità e dalla memoria disponibile.
